% !TEX program = xelatex
\documentclass[11pt,a4paper]{ctexart}

% ==================== 颜色 ====================
\usepackage{xcolor}
\definecolor{maincolor}{HTML}{1A7A6B}       % 主色（标题、顶栏、框边框）
\definecolor{accent}{HTML}{D4A843}          % 强调色（关键词高亮）
\definecolor{boxbg}{HTML}{F0F7F5}           % 问题框背景
\definecolor{borderblue}{HTML}{2E86AB}      % 方法框边框
\definecolor{proofbg}{HTML}{FDF8F0}         % 方法框背景
\definecolor{summarybg}{HTML}{F5F0FF}       % 总结框背景

% ==================== 页面 ====================
\usepackage[top=2.2cm, bottom=2.2cm, left=2.5cm, right=2.5cm, headheight=23.12pt]{geometry}

% ==================== 字体 ====================
\setmainfont{TeX Gyre Pagella}

% ==================== 数学 ====================
\usepackage{amsmath,amssymb,amsfonts,mathtools}

% ==================== 彩色框 ====================
\usepackage{tcolorbox}
\tcbuselibrary{skins,breakable,most}

% ==================== 页眉页脚 ====================
\usepackage{fancyhdr}
\pagestyle{fancy}
\fancyhf{}
\fancyhead[L]{\color{maincolor}\sffamily\bfseries 数学讲义}
\fancyhead[R]{\color{gray}\sffamily \today}
\fancyfoot[C]{\color{gray}\sffamily\thepage}
\renewcommand{\headrulewidth}{1.5pt}
\renewcommand{\headrule}{\hbox to\headwidth{\color{maincolor}\leaders\hrule height \headrulewidth\hfill}}

% ==================== 章节标题样式 ====================
\usepackage{titlesec}
\titleformat{\section}
  {\color{maincolor}\sffamily\LARGE\bfseries}
  {\thesection}{1em}{}
  [\vspace{0.3cm}{\color{maincolor!40}\titlerule[2pt]}]

\titleformat{\subsection}
  {\color{maincolor!80}\sffamily\Large\bfseries}
  {\thesubsection}{1em}{}

% ==================== 表格 ====================
\usepackage{array,booktabs,colortbl}

% ==================== 杂项 ====================
\usepackage{enumitem}
\setlist{nosep, leftmargin=*}
\usepackage[hidelinks]{hyperref}

% ==================== 自定义命令 ====================
\newcommand{\highlight}[1]{\colorbox{accent!20}{\sffamily\bfseries #1}}

% ==================== 自定义环境 ====================

% 1. 问题框 —— 青绿边框 + 浅绿背景
\newtcolorbox[auto counter]{problembox}[1][]{
  enhanced,
  colframe=maincolor,
  colback=boxbg,
  coltitle=white,
  fonttitle=\sffamily\bfseries,
  title={\large 问题},
  attach boxed title to top left={xshift=4mm, yshift=-2mm},
  boxed title style={colback=maincolor, sharp corners},
  sharp corners,
  boxrule=1.2pt,
  left=4mm, right=4mm, top=5mm, bottom=4mm,
  before skip=1.2em, after skip=1.2em,
  #1
}

% 2. 方法/证明框 —— 蓝边框 + 暖白背景
\newtcolorbox[auto counter]{methodbox}[2][]{
  enhanced,
  colframe=borderblue,
  colback=proofbg,
  coltitle=white,
  fonttitle=\sffamily\bfseries,
  title={#2},
  attach boxed title to top left={xshift=4mm, yshift=-2mm},
  boxed title style={colback=borderblue, sharp corners},
  sharp corners,
  boxrule=1.2pt,
  left=4mm, right=4mm, top=5mm, bottom=4mm,
  before skip=1.2em, after skip=1.2em,
  #1
}

% 3. 总结框 —— 淡紫背景
\newtcolorbox{summarybox}[1][]{
  enhanced,
  colframe=maincolor!60,
  colback=summarybg,
  coltitle=white,
  fonttitle=\sffamily\bfseries,
  title={\large $\blacktriangleright$ 总结},
  attach boxed title to top left={xshift=4mm, yshift=-2mm},
  boxed title style={colback=maincolor!70, sharp corners},
  sharp corners,
  boxrule=1.2pt,
  left=4mm, right=4mm, top=5mm, bottom=4mm,
  before skip=1.5em, after skip=1.2em,
  #1
}

% ==================== 文档开始 ====================
\begin{document}

% ---- 彩色顶栏标题 ----
\begin{tcolorbox}[
  enhanced,
  colframe=maincolor,
  colback=maincolor,
  sharp corners,
  boxrule=0pt,
  height=2.8cm,
  valign=center,
  before skip=-2.2cm,
  after skip=1.5cm,
  grow to left by=2.5cm,
  grow to right by=2.5cm,
]
  \begin{center}
    {\color{white}\sffamily\bfseries\Huge 最小正整数问题}\\[3pt]
    {\color{white!85}\sffamily\large 求具有 $21$ 个正因数的最小正整数}
  \end{center}
\end{tcolorbox}

\vspace{0.5cm}

\begin{flushright}
  {\color{gray}\sffamily \today}
\end{flushright}

% ==================== 正文 ====================

\section{问题陈述}

\begin{problembox}
求具有 $21$ 个正因数的\highlight{最小正整数}。
\end{problembox}

这类问题的关键在于把“正因数个数”转化为素因数分解中指数的乘积，再通过指数分配比较不同结构下的最小值。

\section{约数个数定理}

\begin{methodbox}{核心思路}
先写出正整数 $N$ 的标准素因数分解，再用\highlight{约数个数公式}把题目条件 $d(N)=21$ 转化为指数条件。
\end{methodbox}

\subsection*{标准素因数分解}

设正整数 $N$ 的标准素因数分解为：
\[
  N = p_1^{a_1} \cdot p_2^{a_2} \cdots p_k^{a_k}
\]

其中 $p_1,p_2,\ldots,p_k$ 为不同素数，$a_1,a_2,\ldots,a_k$ 为正整数。

\subsection*{正因数个数公式}

其正因数个数 $d(N)$ 的计算公式为：
\[
  d(N) = (a_1 + 1)(a_2 + 1) \cdots (a_k + 1)
\]

题目要求：
\[
  d(N)=21
\]

因此需要把 $21$ 分解为若干个大于 $1$ 的整数乘积，再反推出对应的指数。

\section{分解目标值}

\begin{methodbox}{核心思路}
将 $21$ 写成大于 $1$ 的整数乘积。每个乘积因子都对应某个指数加一，即 $a_i+1$。
\end{methodbox}

\subsection*{目标值分解}

将数字 $21$ 分解为大于 $1$ 的整数乘积，只有两类情况：
\[
  21 = 21
\]
\[
  21 = 7 \times 3
\]

\subsection*{对应的素因子结构}

由 $d(N)=(a_1+1)(a_2+1)\cdots(a_k+1)$ 可知，对应的 $N$ 的素因子结构必然为以下两种形态：
\[
  N = p_1^{20}
\]
\[
  N = p_1^6 \cdot p_2^2
\]

第一种来自 $a_1+1=21$，所以 $a_1=20$；第二种来自 $a_1+1=7$、$a_2+1=3$，所以指数为 $6$ 和 $2$。

% ==================== 新页 ====================
\newpage

\section{最小值分析}

\begin{methodbox}{核心思路}
为使 $N$ 尽可能小，应将\highlight{较大的指数}分配给\highlight{较小的素数}。因此优先使用最小素数 $2,3,5,\ldots$，并让较大指数落在 $2$ 上。
\end{methodbox}

\subsection*{素数选择原则}

正整数从小到大的素数依次为：
\[
  2,3,5,\dots
\]

在指数结构确定后，要得到最小的 $N$，应从最小的素数开始取值，并让较大的指数对应较小的素数。

\subsection*{单素因子结构}

对于形态一
\[
  N = p_1^{20}
\]

取最小素数 $2$，得到：
\[
  N = 2^{20} = 1048576
\]

\subsection*{双素因子结构}

对于形态二
\[
  N = p_1^6 \cdot p_2^2
\]

将指数 $6$ 和 $2$ 分别分配给素数 $2$ 和 $3$：
\[
  N = 2^6 \cdot 3^2
\]

计算得：
\[
  N = 64 \times 9 = 576
\]

\subsection*{两种结构比较}

对比两种形态取得的极小值：
\[
  576 < 1048576
\]

因此满足条件的最小正整数来自第二种结构。

\section{总结}

\begin{summarybox}

\subsection*{推导结构}

\begin{center}
\begin{tabular}{@{}>{\sffamily\bfseries}p{3.2cm} p{10.8cm}@{}}
\toprule
\textbf{步骤} & \textbf{内容} \\
\midrule
约数个数定理 & 若 $N=p_1^{a_1}p_2^{a_2}\cdots p_k^{a_k}$，则 $d(N)=(a_1+1)(a_2+1)\cdots(a_k+1)$ \\
目标条件 & 题目要求 $d(N)=21$ \\
目标分解 & $21=21$ 或 $21=7\times3$ \\
结构一 & $N=p_1^{20}$，最小取 $2^{20}=1048576$ \\
结构二 & $N=p_1^6\cdot p_2^2$，最小取 $2^6\cdot3^2=576$ \\
比较结论 & $576<1048576$ \\
\bottomrule
\end{tabular}
\end{center}

\vspace{1em}

\subsection*{结论}

满足条件的最小正整数为：

\begin{center}
\begin{tcolorbox}[
  enhanced, colframe=accent!80, colback=accent!10,
  sharp corners, boxrule=1.5pt, width=6cm, center,
]
  {\color{maincolor}\sffamily\bfseries\LARGE $576$}
\end{tcolorbox}
\end{center}

验证：
\[
  576 = 2^6 \cdot 3^2
\]
所以
\[
  d(576)=(6+1)(2+1)=7\times3=21
\]

\end{summarybox}

\end{document}
